% ==============================================================================
% Programmatically Generated Manuscript Tables for Cardiorespiratory Conformal Triage
% Source Data: PhysioNet BIDMC (53 ICU Patients) | 5-Fold Patient-Disjoint Evaluation
% Generated by generate_paper_tables.py directly from verified CSV outputs.
% ==============================================================================

\begin{table}[htbp]
\centering
\caption{Predictive accuracy comparison between Extreme Gradient Boosting (XGBoost) and the central-tendency trivial baseline (training median) under 5-fold patient-disjoint cross-validation repeated across 5 independent fold partitions (25 runs total). Errors reported as Mean Absolute Error (MAE) in breaths per minute (bpm). 95\% confidence intervals (CI) derived from 1,000 patient-clustered bootstrap resamples.}
\label{tab:predictive_accuracy}
\begin{tabular}{lccc}
\toprule
\textbf{Model} & \textbf{Mean Test MAE (bpm)} & \textbf{Between-Repeat Std (bpm)} & \textbf{Per-Repeat MAEs (bpm)} \\
\midrule
\textbf{XGBoost (Hist, CPU)} & \textbf{1.81} & 0.09 & [1.86, 1.92, 1.87, 1.68, 1.73] \\
Trivial Baseline (Train Median) & 2.13 & 0.03 & [2.11, 2.19, 2.12, 2.10, 2.11] \\
\bottomrule
\end{tabular}
\vspace{1mm}\\
\footnotesize{\textit{Note}: Paired patient-clustered bootstrap difference ($\Delta = \text{XGB} - \text{Median}$): $-0.31$~bpm, 95\% CI: $[-0.51, -0.11]$~bpm (zero excluded, $p < 0.05$).}
\end{table}


\begin{table*}[htbp]
\centering
\caption{Empirical coverage and efficiency comparison across conformal score functions and cross-conformal Jackknife+/CV+ at 90.0\% nominal target confidence under repeated patient-disjoint cross-validation (25 folds: 5 repeat seeds $\times$ 5 folds). Test MAE and prediction interval widths reported in breaths per minute (bpm). Coverage reported in percentage (\%). AUROC evaluates discriminative capacity to flag estimation errors exceeding 3.0~bpm on clean test windows.}
\label{tab:conformal_scores}
\begin{tabular}{lcccccc}
\toprule
\textbf{Conformal Method} & \textbf{Point MAE (bpm)} & \textbf{Mean Cov (\%)} & \textbf{Cov Spread [Min, Max]} & \textbf{Mean Width (bpm)} & \textbf{AUROC ($|\text{Err}| > 3$~bpm)} & \textbf{Patient Cov (\%)} \\
\midrule
Trivial Baseline (Train Median) & 2.13 $\pm$ 0.41 & 86.6\% $\pm$ 10.4\% & 32.7\% [67.3\%, 100.0\%] & 9.59 $\pm$ 2.84 & 0.50 (constant) & 85.3\% $\pm$ 26.0\% \\
XGBoost (Absolute Residual) & 1.81 $\pm$ 0.43 & 87.5\% $\pm$ 9.8\% & 34.5\% [65.5\%, 100.0\%] & 8.54 $\pm$ 2.86 & 0.50 (constant) & 86.3\% $\pm$ 22.6\% \\
\textbf{XGBoost (Difficulty-Normalised)} & 1.81 $\pm$ 0.43 & 87.4\% $\pm$ 9.7\% & \textbf{31.8\%} [68.2\%, 100.0\%] & 9.30 $\pm$ 4.87 & 0.68 $\pm$ 0.12 & 86.3\% $\pm$ 20.5\% \\
XGBoost (Mondrian Terciles) & 1.81 $\pm$ 0.43 & 84.5\% $\pm$ 10.3\% & 38.0\% [62.0\%, 100.0\%] & 7.71 $\pm$ 2.56 & 0.52 $\pm$ 0.12 & 83.3\% $\pm$ 23.7\% \\
XGBoost (Spectral Sharpness) & 1.81 $\pm$ 0.43 & 87.4\% $\pm$ 10.5\% & 35.5\% [64.5\%, 100.0\%] & 8.64 $\pm$ 3.16 & 0.60 $\pm$ 0.08 & 86.2\% $\pm$ 22.8\% \\
\textbf{Jackknife+ / CV+ (Cross-Fitting)} & 1.73 $\pm$ 0.38 & 90.7\% $\pm$ 7.9\% & \textbf{29.7\%} [70.2\%, 100.0\%] & 8.98 $\pm$ 1.16 & 0.57 $\pm$ 0.16 & 89.4\% $\pm$ 20.6\% \\
\bottomrule
\end{tabular}
\vspace{1mm}\\
\footnotesize{\textit{Note}: Bootstrap 95\% CI on fold-to-fold coverage SD difference (CV+ $-$ Normalised): $[-4.7\%, +1.4\%]$ (includes zero; 1,000 fold-level resamples).}
\end{table*}


\begin{table}[htbp]
\centering
\caption{Discriminative capacity (AUROC) for detecting time-domain photoplethysmogram (PPG) corruptions across mild, moderate, and severe disturbance tiers. Evaluates raw conformal prediction interval width, fold-normalized interval width, and the in-band Signal Quality Index (SQI) baseline. Signal-to-noise ratio (SNR) reported in decibels (dB). AUROC is dimensionless in [0, 1]. 95\% confidence intervals (CI) computed using 1,000 patient-clustered bootstrap resamples.}
\label{tab:corruption_auroc}
\begin{tabular}{lcccc}
\toprule
\textbf{Corruption Severity} & \textbf{SNR (dB)} & \textbf{Raw Width AUROC [95\% CI]} & \textbf{Fold-Norm Width AUROC [95\% CI]} & \textbf{In-Band SQI AUROC [95\% CI]} \\
\midrule
\textbf{Mild Artifact} & +5.0~dB & 0.63 [0.57, 0.69] & 0.66 [0.59, 0.73] & \textbf{0.90 [0.88, 0.92]} \\
\textbf{Moderate Artifact} & -5.5~dB & 0.67 [0.59, 0.74] & 0.69 [0.62, 0.77] & \textbf{0.99 [0.99, 1.00]} \\
\textbf{Severe Artifact} & -12.6~dB & 0.73 [0.65, 0.81] & 0.75 [0.67, 0.82] & \textbf{1.00 [1.00, 1.00]} \\
\bottomrule
\end{tabular}
\end{table}


\begin{table}[htbp]
\centering
\caption{Discriminative capacity (AUROC) for detecting respiratory rate estimation errors exceeding 3.0~bpm on uncorrupted clean PPG windows. Evaluates unnormalized conformal interval width, fold-normalized width, the signal quality index (SQI), and their multi-domain logistic triage combination. AUROC is dimensionless in [0, 1]. Error threshold in breaths per minute (bpm). 95\% confidence intervals (CI) derived from 1,000 patient-clustered bootstrap resamples.}
\label{tab:clean_error_detection_auroc}
\begin{tabular}{lcc}
\toprule
\textbf{Predictor Rule} & \textbf{Pooled AUROC [95\% CI]} & \textbf{Fold AUROC (Mean $\pm$ Std)} \\
\midrule
Raw Interval Width & 0.61 [0.50, 0.72] & 0.68 $\pm$ 0.15 \\
Fold-Normalized Width & 0.70 [0.60, 0.80] & 0.68 $\pm$ 0.15 \\
Signal Quality Index (SQI) & 0.60 [0.51, 0.70] & 0.63 $\pm$ 0.14 \\
\textbf{Multi-Domain Combination (Width + SQI)} & \textbf{0.75 [0.63, 0.86]} & 0.71 $\pm$ 0.11 \\
\bottomrule
\end{tabular}
\end{table}


\begin{table*}[htbp]
\centering
\caption{Selective prediction risk-coverage operating points on clean test windows across retention fractions from 40\% through 100\%. Retention reported in percentage (\%). Errors reported as Mean Absolute Error (MAE) in breaths per minute (bpm). Compares raw conformal width, fold-normalized width, multi-domain combination score, oracle lower bound, and random abstention. 95\% confidence intervals (CI) derived from 1,000 patient-clustered bootstrap resamples.}
\label{tab:risk_coverage_retention}
\begin{tabular}{cccccc}
\toprule
\textbf{Retention} & \textbf{Raw Width MAE (bpm)} & \textbf{Fold-Norm Width MAE (bpm)} & \textbf{Combination Score MAE (bpm)} & \textbf{Oracle Lower Bound (bpm)} & \textbf{Random Abstention (bpm)} \\
\midrule
\textbf{40\%} & 1.20 [0.85, 1.93] & \textbf{1.01 [0.74, 1.52]} & 1.14 [0.79, 1.59] & 0.33 [0.26, 0.44] & 1.82 [1.36, 2.36] \\
\textbf{50\%} & 1.59 [1.03, 2.21] & 1.26 [0.84, 1.80] & \textbf{1.15 [0.84, 1.57]} & 0.43 [0.34, 0.56] & 1.82 [1.35, 2.35] \\
\textbf{60\%} & 1.71 [1.18, 2.29] & 1.41 [0.96, 1.99] & \textbf{1.18 [0.86, 1.59]} & 0.54 [0.43, 0.71] & 1.78 [1.32, 2.29] \\
\textbf{70\%} & 1.74 [1.26, 2.34] & 1.53 [1.08, 2.08] & \textbf{1.27 [0.91, 1.80]} & 0.68 [0.54, 0.90] & 1.79 [1.34, 2.31] \\
\textbf{80\%} & 1.83 [1.35, 2.39] & 1.64 [1.20, 2.16] & \textbf{1.41 [1.01, 1.91]} & 0.87 [0.68, 1.19] & 1.79 [1.33, 2.32] \\
\textbf{85\%} & 1.90 [1.38, 2.44] & 1.66 [1.23, 2.16] & \textbf{1.49 [1.09, 1.98]} & 1.00 [0.78, 1.42] & 1.80 [1.34, 2.33] \\
\textbf{90\%} & 1.85 [1.36, 2.43] & 1.67 [1.26, 2.22] & \textbf{1.57 [1.17, 2.06]} & 1.19 [0.89, 1.69] & 1.79 [1.34, 2.33] \\
\textbf{95\%} & 1.84 [1.35, 2.39] & 1.77 [1.31, 2.33] & \textbf{1.64 [1.26, 2.22]} & 1.46 [1.06, 1.99] & 1.80 [1.34, 2.33] \\
\textbf{100\%} & 1.80 [1.34, 2.33] & \textbf{1.80 [1.34, 2.33]} & 1.80 [1.34, 2.33] & 1.80 [1.34, 2.33] & 1.80 [1.34, 2.33] \\
\bottomrule
\end{tabular}
\end{table*}


\begin{table}[htbp]
\centering
\caption{External out-of-distribution validation on CapnoBase benchmark stratified by population cohort. The adult-only cohort ($N=13$) represents the primary, domain-appropriate generalization test (both BIDMC and CapnoBase adults within the 6--36~bpm designed passband). Pediatric cases ($N=29$) exhibit a physiological frequency mismatch (4.48\% of windows $> 36$~bpm), causing passband signal attenuation. Point MAE and prediction interval widths reported in breaths per minute (bpm). Coverage reported in percentage (\%). All values reported as mean $\pm$ standard deviation across 5 BIDMC fold models. AUROC evaluates clean-window error discrimination ($|\text{Error}| > 3.0$~bpm).}
\label{tab:external_validation_capnobase}
\begin{tabular}{lcccc}
\toprule
\textbf{Metric} & \textbf{In-Domain (BIDMC ICU)} & \textbf{External Adult (Primary)} & \textbf{External Pediatric} & \textbf{Full Cohort (All 42)} \\
\midrule
Point Prediction MAE (bpm) & 1.78 $\pm$ 0.55 & 6.21 $\pm$ 0.36 & 6.39 $\pm$ 0.29 & 6.33 $\pm$ 0.31 \\
Empirical Coverage (90\% nominal target) & 88.1\% $\pm$ 11.9\% & 45.3\% $\pm$ 18.4\% & 59.8\% $\pm$ 21.6\% & 55.3\% $\pm$ 20.3\% \\
Mean Prediction Interval Width (bpm) & 10.72 $\pm$ 4.87 & 11.94 $\pm$ 5.39 & 15.32 $\pm$ 6.86 & 14.29 $\pm$ 6.32 \\
Patient-Level Coverage (\%) & 87.1\% $\pm$ 19.2\% & 44.5\% $\pm$ 37.9\% & 59.8\% $\pm$ 37.4\% & 55.1\% $\pm$ 38.5\% \\
Error Detection AUROC ($|\text{Error}| > 3.0$~bpm) & 0.712 $\pm$ 0.127 & 0.528 $\pm$ 0.148 & 0.619 $\pm$ 0.100 & 0.577 $\pm$ 0.101 \\
Selective Triage MAE (80\% Retention, bpm) & -- & 6.17 $\pm$ 0.44 & 5.97 $\pm$ 0.35 & 5.98 $\pm$ 0.31 \\
\bottomrule
\end{tabular}
\vspace{1mm}\\
\footnotesize{\textit{Note}: Evaluated zero-shot using frozen BIDMC-trained models and conformal calibrators. The adult-only cohort ($N=13$, 2,689 windows; age $\ge 18$) is the primary domain generalization benchmark (both cohorts within the 6--36~bpm designed filter passband). Pediatric cases ($N=29$, 6,075 windows; age $< 18$) incorporate a direct frequency-range mismatch (4.48\% of windows $> 36$~bpm, exceeding the 0.10--0.60~Hz Butterworth passband ceiling). Conformal interval width expands adaptively (Adult: 11.94~bpm, Pediatric: 15.32~bpm), and selective triage (80\% retention) reduces error across all cohorts.}\par
\end{table}
